\chapter{Полное радиальное уравнение дефекта и его вторая вариация}

Предыдущая проверка доказала, что абсолютные единицы длины и энергии не
следуют из текущего родителя. Это не мешает поставить безразмерный вопрос:
существует ли вообще стационарный дефект полного функционала, или конечная
энергия прежнего пробного профиля была лишь свойством удачного анзаца.

\section{Литературная и проектная граница}

В моделях Скирма радиальный анзац обычно приводит к нелинейной краевой
задаче, которую решают непосредственно и лишь затем исследуют вторую
вариацию; пример такого анализа дан в \path{arXiv:hep-ph/0106150}.
Математическая теория функционалов Фаддеева--Скирма связывает
четырёхпроизводный член с существованием топологически устойчивых
минимизаторов, но подчёркивает зависимость строгих теорем от области и
функционального класса; см. \path{arXiv:0907.0899} и
\path{arXiv:math-ph/0608042}. Радиальная положительность не заменяет полный
спектральный анализ: даже в стандартной модели Скирма симметричные и
несимметричные возмущения являются разными задачами; см.
\path{arXiv:0712.1510}.

Поэтому настоящий результат имеет точную область действия: он проверяет
сферически-симметричный сектор составного проекторного ежа на
$\mathbb R^3$, но не объявляет доказанной устойчивость относительно всех
пяти компонент поля $Q$.

\section{Точная радиальная редукция}

Положим
\begin{equation}
 Q=\Delta u(r)\left(P-\frac{I_3}{3}\right),
 \qquad
 P=\widehat x\widehat x^T,
 \qquad
 A_Q=u(r)^2[P,dP].
\end{equation}
Граничные условия гладкого единичного ежа имеют вид
\begin{equation}
 u(0)=0,
 \qquad
 u(\infty)=1.
 \label{eq:v6-bosonic-radial-boundary-conditions}
\end{equation}

Прямое матричное вычисление даёт
\begin{align}
 |D_QQ|^2
 &=\frac23\Delta^2u'^2
 +\frac{4\Delta^2u^2(1-u^2)^2}{r^2},\\
 |F_{A_Q}|^2
 &=\frac{16u^2u'^2}{r^2}
 +\frac{2u^4(u^2-2)^2}{r^4},\\
 V(u)&=\Delta^4(1-u^2)^2.
 \label{eq:v6-bosonic-radial-densities}
\end{align}
Максимальное расхождение этих формул с независимым трёхмерным
декартовым вычислением равно $1.9\cdot10^{-11}$.

Для положительных коэффициентов $c_D,c_F,c_V$ энергия без общего множителя
$4\pi$ равна
\begin{equation}
 \mathcal E[u]=\int_0^\infty
 \left(A(r,u)u'^2+W(r,u)\right)dr,
 \label{eq:v6-bosonic-radial-functional}
\end{equation}
где
\begin{align}
 A(r,u)&=\frac23c_D\Delta^2r^2+16c_Fu^2,\\
 W(r,u)&=4c_D\Delta^2u^2(1-u^2)^2
 +\frac{2c_Fu^4(u^2-2)^2}{r^2}
 +c_V\Delta^4r^2(1-u^2)^2.
\end{align}
Уравнение Эйлера--Лагранжа принимает компактный вид
\begin{equation}
 2A u''+2A_r u'+A_u u'^2-W_u=0.
 \label{eq:v6-bosonic-radial-euler-lagrange}
\end{equation}

\section{Численное решение полного уравнения}

Краевая задача решена на $10^{-4}\le r\le50$ с аналитической поправкой
кривизненного хвоста
\begin{equation}
 E_F(50,\infty)=\frac{8\pi c_F}{50}.
\end{equation}
При единичных коэффициентах получены
\begin{align}
 r_{1/2}&=1.0394444805\ldots,\\
 E_D&=11.0306107926\ldots,\\
 E_F&=35.3958853096\ldots,\\
 E_V&=8.1217582001\ldots,\\
 E_{\rm exact}&=54.5482543023\ldots.
 \label{eq:v6-bosonic-exact-radial-energy}
\end{align}
Здесь $r_{1/2}$ определяется условием $u(r_{1/2})=1/2$.

Лучший профиль старого однопараметрического семейства имел энергию
\begin{equation}
 E_{\rm old}=60.7698954459\ldots.
\end{equation}
Полное решение ниже него на $10.2380\ldots\%$.
Следовательно, стационарная конфигурация не является артефактом выбранной
рациональной формы $r^2/(r^2+R^2)$.

Вириальное тождество Деррика
\begin{equation}
 E_D-E_F+3E_V=0
 \label{eq:v6-bosonic-radial-virial-identity}
\end{equation}
выполнено с остатком $8.4\cdot10^{-8}$ после поправки хвоста.

\begin{theorem}[Существование численного радиального решения]
Для $c_D=c_F=c_V=1$ краевая задача
\eqref{eq:v6-bosonic-radial-boundary-conditions}--
\eqref{eq:v6-bosonic-radial-euler-lagrange} имеет гладкое монотонное
численное решение с конечной энергией. Его энергия строго ниже минимума
прежнего масштабированного пробного семейства.
\end{theorem}

\section{Окно положительных коэффициентов}

Чтобы не принять случайную точку за механизм, вычисление повторено для
семи наборов:
\begin{equation}
 (c_D,c_F,c_V)\in
 \left\{
 (1,1,1),
 (1/4,1,1),(4,1,1),
 (1,1/4,1),(1,4,1),
 (1,1,1/4),(1,1,4)
 \right\}.
\end{equation}
Во всех случаях решатель сошёлся с максимальным среднеквадратичным
остатком менее $2.0\cdot10^{-7}$, профиль оказался монотонным, точная
энергия была ниже соответствующего масштабированного пробного семейства,
а вириальный остаток не превысил $3.0\cdot10^{-7}$.

Это конечное вычислительное окно, а не теорема для всего положительного
октанта коэффициентов. Тем не менее оно показывает, что решение не требует
тонкой настройки отношения трёх секторов.

\section{Радиальная вторая вариация}

Для возмущения $u\mapsto u+\varepsilon\eta$ с
$\eta(0)=\eta(\infty)=0$ квадратичная форма приводится к оператору
Штурма--Лиувилля
\begin{equation}
 \mathcal H_{\rm rad}\eta
 =-\frac{d}{dr}\left(2A\frac{d\eta}{dr}\right)
 +U_{\rm rad}(r)\eta,
 \label{eq:v6-bosonic-radial-hessian}
\end{equation}
где
\begin{equation}
 U_{\rm rad}
 =A_{uu}u'^2+W_{uu}
 -2\frac{d}{dr}(A_u u').
\end{equation}

Число отрицательных собственных значений вычислялось не только прямой
диагонализацией, но и устойчивым штурмовским подсчётом знаков главных
поворотов трёхдиагональной матрицы. Во всех семи случаях отрицательных
радиальных мод не найдено.

\begin{proposition}[Радиальная устойчивость в проверенном окне]
Полученные семь стационарных решений являются локальными минимумами внутри
сферически-симметричного сектора: дискретизированный оператор
\eqref{eq:v6-bosonic-radial-hessian} не имеет отрицательных собственных
значений.
\end{proposition}

\section{Граница результата}

Поле $Q$ имеет пять локальных компонент. Радиальный анзац проверяет только
амплитудное направление, совместимое с диагональным ежом. Он не охватывает:
\begin{enumerate}
 \item двухосные возмущения ядра;
 \item угловые гармоники с $\ell\ge1$;
 \item смешанные вариации $Q$ и составной связности вне условия
 $A_Q=[Q,dQ]/\Delta^2$;
 \item распад на несферические дефектные фрагменты;
 \item временные возмущения до фиксации общей лоренцевой нормы.
\end{enumerate}

\begin{nogocriterion}[Радиальная устойчивость не равна полной]
Отсутствие отрицательной моды в операторе
$\mathcal H_{\rm rad}$ не разрешает называть дефект полностью устойчивой
частицей. Следующая проверка обязана разложить все пять компонент $Q$ по
сферическим гармоникам и проверить по меньшей мере секторы $\ell=0,1,2$.
\end{nogocriterion}

\begin{protocol}
\begin{enumerate}
 \item точная радиальная редукция: \textbf{получена};
 \item декартовый остаток редукции: \textbf{$<2\cdot10^{-11}$};
 \item полное радиальное уравнение: \textbf{решено численно};
 \item единичная энергия: \textbf{$54.5482543023\ldots$};
 \item выигрыш относительно пробного семейства: \textbf{$10.2380\ldots\%$};
 \item вириальное тождество: \textbf{выполнено};
 \item семь положительных наборов коэффициентов: \textbf{сошлись};
 \item радиальная отрицательная мода: \textbf{не найдена};
 \item полная нерадиальная устойчивость: \textbf{не доказана};
 \item следующая проверка: \textbf{нерадиальный спектр пятикомпонентного поля};
 \item устойчивая наблюдаемая частица: \textbf{ещё не получена}.
\end{enumerate}
\end{protocol}

Машинный сертификат сохранён в
\path{s2t_v6_bosonic_defect_full_euler_lagrange_stability_gate_results.json}.